\documentclass[10pt]{article}
\usepackage[T1]{fontenc}
\usepackage[utf8]{inputenc}
\usepackage[french]{babel}
\usepackage{amsmath,amssymb,mathrsfs,mathpazo}
\usepackage{geometry}
\geometry{a4paper,margin=2cm}
\usepackage{graphicx,array,tabularx,tabularray,colortbl,multirow,multicol,booktabs,enumitem,pifont,ulem,fancybox,diagbox,cancel}
\usepackage{pstricks,pstricks-add,pst-plot,pst-tree,pst-node,pst-text,pst-func,pst-math,pst-bspline,pst-3dplot,pst-eucl,pst-solides3d}
\let\pscalculate\undefined
\usepackage{pst-bezier}
\usepackage{tikz,pgfplots,tkz-tab}
\usetikzlibrary{arrows,arrows.meta,patterns,calc,positioning,trees,decorations.markings,shapes.geometric,shapes.misc,shapes.arrows,intersections,decorations.pathmorphing,decorations.pathreplacing,backgrounds,shadows,babel}
\pgfplotsset{compat=1.18}
\usepgfplotslibrary{fillbetween}
\usepackage[np]{numprint}
\usepackage{listings,makecell,pas-tableur,twoopt}
\let\barre\undefined
\usepackage[pstricks]{tablvar}
\definecolor{darkgreen}{RGB}{0,100,0}
\providecommand{\red}{\color{red}}
\providecommand{\blue}{\color{blue}}
\providecommand{\green}{\color{green}}
\providecommand{\black}{\color{black}}
\providecommand{\white}{\color{white}}
\providecommand{\gray}{\color{gray}}
\providecommand{\cyan}{\color{cyan}}
\providecommand{\magenta}{\color{magenta}}
\providecommand{\yellow}{\color{yellow}}
\providecommand{\violet}{\color{violet}}
\providecommand{\orange}{\color{orange}}
\providecommand{\euro}{\texteuro}
\providecommand{\EUR}{\texteuro}
\providecommand{\ieme}{\textsuperscript{e}}
\providecommand{\degres}{\ensuremath{^\circ}}
\providecommand{\degree}{\ensuremath{^\circ}}
\providecommand{\vv}[1]{\overrightarrow{#1}}
\providecommand{\Coord}[1]{#1}
\providecommand{\decosix}{}
\providecommand{\starredbullet}{\ensuremath{\bullet}}
\pagestyle{empty}
\setlength{\parindent}{0pt}
\usepackage{adjustbox}\DeclareUnicodeCharacter{2236}{:}\DeclareUnicodeCharacter{2190}{\ensuremath{\leftarrow}}
\providecommand{\R}{\mathbb{R}}
\providecommand{\N}{\mathbb{N}}
\providecommand{\Z}{\mathbb{Z}}
\providecommand{\Q}{\mathbb{Q}}
\providecommand{\D}{\mathbb{D}}
\providecommand{\C}{\mathbb{C}}
\providecommand{\e}{{\mathrm{e}}}
\providecommand{\vect}[1]{\overrightarrow{#1}}
\providecommand{\vectt}[1]{\overrightarrow{\mathrm{#1}}}
\providecommand{\vvt}[1]{\overrightarrow{\mathrm{#1}}}
\providecommand{\np}[2][]{{#2~#1}}
\providecommand{\vv}[1]{{\overrightarrow{#1}}}
\renewcommand{\d}{{\mathrm{d}}}
\providecommand{\degres}{{^\circ}}
\providecommand{\degre}{{^\circ}}
\providecommand{\blue}{\color{blue}}
\providecommand{\red}{\color{red}}
\providecommand{\enskip}{\quad}
\providecommand{\llbracket}{\lbrack\!\lbrack}
\providecommand{\rrbracket}{\rbrack\!\rbrack}
\providecommand{\euro}{{\text{€}}}
\providecommand{\no}{{\text{n°}}}
\providecommand{\interval}[2]{{[#1;#2]}}
\providecommand{\intervFO}[2]{{[#1;#2[}}
\providecommand{\dotfill}{\ldots}

\providecommand{\ds}{\displaystyle}
\providecommand{\textding}[1]{#1}
\providecommand{\Oij}{\(\mathrm{O};\vec{\imath},\vec{\jmath}\)}
\providecommand{\Oijk}{\(\mathrm{O};\vec{\imath},\vec{\jmath},\vec{k}\)}
\providecommand{\Ouv}{\(\mathrm{O};\vec{u},\vec{v}\)}
\providecommand{\intervFF}[2]{\ensuremath{[#1;#2]}}
\renewcommand{\euro}{\ifmmode\text{\texteuro}\else\texteuro\fi}
\DeclareUnicodeCharacter{00B0}{\ifmmode^{\circ}\else\textdegree\fi}
\DeclareUnicodeCharacter{20AC}{\ifmmode\text{\texteuro}\else\texteuro\fi}
% Énoncés : Éducation nationale. Transcription et corrigés : APMEP, Denis Vergès et contributeurs. Découpage et classement : CoopMaths, Sébastien Lozano, Éric Elter et contributeurs. Adaptation de la présentation : Maths Explorer.
% CC BY-SA : https://coopmaths.fr/www/about/
% https://coopmaths.fr/alea/static/bac/2024/tex/bac_2024_06_sujet2_polynesie_4_cor.tex
\begin{document}
\noindent\textit{APMEP / CoopMaths. CC BY-SA. Présentation adaptée par Maths Explorer.}\par\medskip
\medskip




















\begin{enumerate}
	\item 
On a $\vect{\text{AB}}\begin{pmatrix}5\\-1\\- 13\end{pmatrix}$ et $\vect{\text{AC}}\begin{pmatrix}2\\-2\\- 10\end{pmatrix}$.

Ces vecteurs ne sont pas colinéaires donc les points A, B et C  sont distincts et non alignés : ils définissent le plan $\mathcal{P} = $ (ABC).
\end{enumerate}

Dans la suite, on considère le vecteur $\vect{n} \begin{pmatrix}2\\-3\\1\end{pmatrix}$.

\begin{enumerate}[resume]
	\item 
	\begin{enumerate}
		\item 
On a $\vect{n} \cdot \vect{\text{AB}} = 10  + 3 - 13 = 0$ et  $\vect{n} \cdot \vect{\text{AB}} = 4  + 6 - 10 = 0$.

Le vecteur $\vect{n}$ orthogonal à deux vecteurs non colinéaires du plan (ABC) est normal à ce plan
		\item 
On sait qu'alors :
		
$M(x~;~y~;~z) \in  (\text{ABC}) \iff 2x - 3y + 1z + d = 0, \: d \in \R$.
		
Par exemple A$(-1~;~-1~;~17)\in (\text{ABC}) \iff - 2 + 3 + 17 + d = 0 \iff d = -18$, donc 

$M(x~;~y~;~z) \in  (\text{ABC}) \iff 2x - 3y + z - 18 = 0$.
\end{enumerate}
	\item 

\[ d~:~\left\{\begin{array}{l c l}
x&=&3t+2\\y&=&\phantom{3}t+5\\z&=&4t+1\end{array}\right.\text{, avec } t \in \R. \]
	\begin{enumerate}
		\item 
On sait que les coordonnées d'un vecteur directeur $\vect{d}$ de $d$ sont les coefficients de $t$, soit $\vect{d}\begin{pmatrix}3\\1\\4\end{pmatrix}$
		\item 
		Les coordonnées de E commun à $d$ et au plan (ABC) vérifient les équations paramétriques de $d$ et l'équation cartésienne de (ABC) soit le système
		
		$\left\{\begin{array}{l c l}
x&=&3t+2\\y&=&\phantom{3}t+5\\z&=&4t+1\\
2x - 3y + 1z - 18 &=& 0
\end{array}\right.\text{, avec } t \in \R.$

En remplaçant $x, \:y$ et $z$ par leurs valeurs en fonction de $t$ dans la dernière équation on obtient :

$2(3t + 2)- 3(t + 5) + 4t + 1 - 18 = 0 \iff 6t + 4 - 3y - 15 + 4t + 1 - 18 = 0\iff $

$7t - 28 = 0 \iff 7t = 28 \iff t = 4$.

D'où les coordonnées de E en remplaçant $t$ par dans les équations de $d$ : E(14~;~9~;~17)
	\end{enumerate}
	\item 
La droite $\Delta$ contient D et a pour vecteur directeur un vecteur normal au plan (ABC) soit le vecteur $\vect{n}$.

Une équation de $\Delta$ est donc $\left\{\begin{array}{l c l}
x&=&2t+2\\y&=&-3t+5\\z&=&t+1\\
2x - 3y + 1z - 18 &=& 0
\end{array}\right.\text{, avec } t \in \R.$

En résolvant le système obtenu en rajoutant l'équation de $\mathcal{P}$ permet d'obtenir $t = 2$, d'où F$(6~;~-1~;~3)$.

Puisque la droite (DF) est perpendiculaire au plan $\mathcal{P}$, la distance DF est la (plus courte) distance du point au plan $\mathcal{P}$.

Or $\vect{\text{DF}}\begin{pmatrix}4\\-6\\2\end{pmatrix}$, d'où $\left\|\vect{\text{DF}}\right\|^2 = \text{DF}^2 = 4^2 + (-6)^2 + 2^2 = 16 + 36 + 4 = 56 = 4 \times 14$.

On a donc DF $ = \sqrt{4 \times 14} = \sqrt 4 \times \sqrt{14} = 2\sqrt{14}$ (en centaines de metres).
	

	\item 
	

La plus petite distance du point de lancer au plan (ABC) est égale à $2\sqrt{14}$ centaines de mètres : il faut donc calculer le temps mis à la vitesse de 18,6~$\text{m.s}^{-1}$ par un drone pour effectuer ce parcours de D à F.

On sait que $v = \dfrac{d}{t}$, ou $t = \dfrac{d}{v} = \dfrac{\text{DF}}{v} = \dfrac{2\sqrt{14} \times 100}{18,6} \approx 40,23$ (secondes).

Conclusion : le drone n'arrivera pas à temps.
\end{enumerate}

\end{document}
